\chapter{タスク定義とシステム構成}

\section{現場で実行する一連のタスク}

机上に複数のプラスチックボトルを散在させる。左作業腕が対象を選択して持ち上げ、右作業腕がキャップを把持して回し、取り外す。左腕はボトル本体を本体用回収箱へ、右腕はキャップをキャップ用回収箱へ投入し、次のボトルへ移る。\ev{E01,E10}

\begin{table}[H]
\centering
\small
\caption{開始、完了、失敗、中止条件}
\begin{tabularx}{\textwidth}{>{\japanesesans\bfseries}p{23mm}X p{43mm}}
\toprule
区分 & 現在の定義 & エビデンスの状態\\
\midrule
開始 & 双腕が実行可能な初期姿勢にあり、未処理ボトルと 3 視点画像が存在 & 実行設定は確認済み。初期配置の固定試験手順は未整備\\
完了 & キャップと本体が分離し、それぞれ異なる回収箱へ入る & 現場で目視確認。自動判定なし\\
典型失敗 & 未把持、逆向き把持、キャップなしで空回し、未分離、誤投入、落下 & 現場記述はあるが発生頻度の記録なし\\
中止 & 操作者停止、装置異常、安全確認不合格 & 収集・制御ソフトに停止と健全性確認機能あり\\
\bottomrule
\end{tabularx}
\end{table}

\section{ハードウェアと双腕の役割}

\begin{table}[H]
\centering
\small
\caption{ハードウェア構成と役割}
\begin{tabularx}{\textwidth}{>{\japanesesans\bfseries}p{28mm}p{44mm}X}
\toprule
構成 & 主な役割 & 確認済み事実と制約\\
\midrule
左作業腕・グリッパ & 本体の選択、把持、持上げ、開栓時の保持、本体投入 & 左 6 関節と 1 グリッパ量を出力。役割はタスク定義と現場観測で確認\\
右作業腕・グリッパ & 口部接近、キャップ把持、回転、分離、キャップ投入 & 連続回転を含む。独立した力覚・触覚記録なし\\
上方カメラ & 作業台全体、ボトル分布、双腕相対位置を観測 & 学習と実機運用で使用。原画像は 640$\times$480\\
左手首カメラ & 左グリッパ、本体、接触部の近接像 & 学習と実機運用で使用\\
右手首カメラ & キャップ、右グリッパ、回転部の近接像 & アテンション重みの平均構成比が最大\\
操作者用示教腕 & 収集担当者が両手で動作を直接提示 & 連続収集、再収集、復帰、複数トリガ方式を支援\\
作業台・回収箱 & ランダムなボトル群と本体/キャップ投入領域 & 寸法、座標、固定配置の公開校正表は未整備\\
\bottomrule
\end{tabularx}
\end{table}

\section{ハードウェアとソフトウェアの閉ループ}

\figfull[.98\textwidth]{software_data_pipeline.pdf}{初年度に構築したデータ―学習―実機運用の閉ループ。専用収集、データ編集、版公開、モデル学習、現場運用、問題のフィードバックを接続した。現在のソフトウェア機能が、過去の全データに同様に適用済みとは限らない。}

収集担当者は専用ソフトウェアで ALOHA を遠隔操作し、ロボット状態、多視点動画、タスク情報を同時保存する。編集センターで閲覧、切出し、ラベル付与、版生成、公開を行う。学習基盤は固定版データを読み込み、モデルと学習履歴を保存する。現場失敗は、学習時間だけを延長するのではなく、データ分類と追加収集へ戻す。

\section{制御周期、状態、行動}

実機システムは 50 Hz、すなわち 20 ms ごとに制御指令を送る。状態と行動はいずれも 14 次元である。
\[
\mathbf{s}_t,\mathbf{a}_t\in\mathbb{R}^{14}
=[q^L_{1:6},g^L,q^R_{1:6},g^R].
\]
$q^L_{1:6},q^R_{1:6}$ は左右作業腕の 6 関節位置、$g^L,g^R$ は左右グリッパ位置である。モデルは一回に $H=50$ 時刻を予測するため、出力形状は $50\times14$ となる。\ev{E06,E07}

\plain{ロボットが毎回「次の 1 秒間をどう動くか」という短い計画を作ると考えるとよい。実行途中で次の計画を用意し、計画切替時の停止を減らす。}

実行時は、動作列の開始後 25 制御周期目から次の動作列を先行計算し、切替条件を満たした時点で新しい動作列へ移る。これは複数予測の単純平均ではない。関節指令にはハードウェア範囲制限があり、グリッパ電流上限も設定されている。\ev{E07}

\limitbox{現行システムは、キャップの分離、本体の箱内投入、タスク成功を自動判定しない。衝突回数、キャップ回転角、接触力も記録していないため、成功条件は現場目視に依存する。}
